%قبل از شروع مسئله ترجیحاً چیزی ننویسید. اما اگر خواستید بنویسید، در قالب لیستی (فقرات و شمارش) بنویسید.
%برای اضافه کردن نمره، بعدی از شروع‌مسئله، [x] را اضافه کنید و به جای x  نمره مورد نظر را قرار دهید. برای زیرمسئله‌ها هم به همین صورت است.
%\شروع‌مسئله[10]
\شروع‌مسئله [8]
% صورت اصلی مسئله را در این‌جا بنویسید.
دو تابع $F$ و $G$ زیر را به کمک درخت بازگشت به صورت مجانبی تحلیل کنید. 
\begin{align*}
	F(0) &= G(0) = 0 \\
	F(n) &= F(\lfloor \frac{2n}{3} \rfloor) + G(n) \\
	G(n) &= F(\lfloor \frac{n}{3} \rfloor) + n
\end{align*}
%اگر مسئله چند قسمت داشت به این صورت عمل کنید:
%\شروع{زیرمسائل}
%\زیرمسئله[0] % مسئلۀ ۱آ
%\زیرمسئله % مسئلۀ ۱ب
%\زیرمسئله % مسئلۀ ۱پ
%\زیرمسئله % مسئلۀ ۱ت
%\پایان{زیرمسائل}
\solution{
	با جایگذاری رابطه دوم در رابطه اول داریم
	\begin{equation*}
		F(n) = F(\lfloor \frac{2n}{3} \rfloor) + F(\lfloor \frac{n}{3} \rfloor) + n 
	\end{equation*}
	برای تحلیل مجانبی تابع بالا می‌توانیم، رابطه بالا را به شکل زیر بازنویسی کنیم
	\begin{equation*}
		F(n) = F(\frac{2n}{3}) + F(\frac{n}{3}) + n
	\end{equation*}
	درخت بازگشت این رابطه را می‌توان به شکل زیر ترسیم کرد.
\شروع{شکل}[h!]
\تنظیم‌ازوسط
  \درج‌تصویر[پهنا=0.5\پهنای‌متن]{src/p1_fig.png}
\شرح{درخت بازگشت $F$}
\برچسب{درخت بازگشت}
\پایان{شکل}	
دقت کنیم که عمق برگ‌های این درخت حداکثر
$\log_{\frac{3}{2}} n$
و حداقل 
$\log_3 n$
است. پس به سادگی می‌توان دید
\begin{equation*}
\underbrace{n + n + \cdots + n}_{\log_3 n} 	\leq F(n) \leq \underbrace{n + n + \cdots + n}_{\log_\frac{3}{2} n}\end{equation*}
به عبارت دیگر
\begin{equation*}
n \log_3 n	\leq F(n) \leq n\log \log_\frac{3}{2} n
\end{equation*}
پس بنا بر رابطه تغییر پایه در لگاریتم‌ها داریم
\begin{equation*}
c_1n \log_2 n	\leq F(n) \leq c_2 \log_2 n
\end{equation*}
که نشان می‌دهد $F(n) \in \Theta(nlgn)$.
همچنین با توجه به رابطه‌ای که برای $G$ داریم
\begin{equation*}
c_1\frac{n}{3} \log_2 \frac{n}{3} + n	\leq G(n) = F(\frac{n}{3}) + n \leq c_1\frac{n}{3} \log_2 \frac{n}{3} + n
\end{equation*}
که نشان می‌دهد
\begin{equation*}
\frac{c_1}{3} n \log_2 n - c_1'n	\leq G(n)  \leq \frac{c_2}{3} n \log_2 n - c_2'n	
\end{equation*}
پس 
\begin{equation*}
	G(n) \in \Theta(nlgn)
\end{equation*}
}




%%%%%%%%%%%%%%%%%%%%%%%%%%%%%%%%%%%%%%%%%%%%%%%%%%%%%
% برخی دستورات مفید برای تمرین‌ها. %
%%%%%%%%%%%%%%%%%%%%%%%%%%%%%%%%%%%%%%%%%%%%%%%%%%%%%

% فرمول‌های نمونه:
%$f(x) = \Theta(x)$
%$T(x, y) \leq \log(x) + 2^y + \binom{2n}{n}$
% {\tt code\_function}


% لیست موارد بدون شماره‌گذاری:
%\شروع{فقرات}
%    \فقره مورد اول 
%        \شروع{فقرات}
%            \فقره مورد اول داخلی 
%                \شروع{فقرات}
%                    \فقره اولین مورد خیلی داخلی 
%                    \فقره دومین مورد خیلی داخلی 
%                \پایان{فقرات}
%            \فقره مورد دوم داخلی 
%        \پایان{فقرات}
%    \فقره مورد دوم 
%\پایان{فقرات}

% لیست موارد شماره‌گذاری شده:
%\شروع{شمارش}
%    \فقره یک مورد 
%    \فقره یک مورد دیگر 
%    \فقره باز هم یک مورد دیگر
%\پایان{شمارش}

% برای نوشتن یک فرمول چند خطی که یک شماره بخوره. برای حالت‌های دیگر به سایت زیر مراجعه کنید:
% https://www.overleaf.com/learn/latex/Aligning_equations_with_amsmath
%\begin{align} 
%    \begin{split}
%        (x+y)^3 &= (x+y)^2(x+y) \\
%                &= (x^2+2xy+y^2)(x+y) \\
%                &= (x^3+2x^2y+xy^2) + (x^2y+2xy^2+y^3) \\
%                &= x^3+3x^2y+3xy^2+y^3
%    \end{split}                                 
%\end{align}

% نوشتن ماتریس و آرایه دوبعدی
%\begin{align}
%    A = 
%    \begin{bmatrix}
%        A_{11} & A_{21} \\
%        A_{21} & A_{22}
%    \end{bmatrix}
%\end{align}

% اضافه کردن تصویر با شرح. برای اضافه کردن تصویر بدون شرح، فقط دستور درج‌تصویر را استفاده کنید.
%\شروع{شکل}[h!]
%\تنظیم‌ازوسط
%  \درج‌تصویر[پهنا=0.5\پهنای‌متن]{img.jpg}
%\شرح{شرح تصویر}
%\برچسب{تصویر:۱}
%\پایان{شکل}
%\پایان{زیرمسائل}

%کد زیر را ببینید.
%%\lstinputlisting{source_filename.py}
%\begin{latin}
%\begin{lstlisting}
%# Import the modules
%import sys
%import random
%
%ans = True
%
%while ans:
%question = raw_input("Ask the magic 8 ball a question: (press enter to quit) ")
%
%answers = random.randint(1,8)
%
%if question == "":
%sys.exit()
%
%elif answers == 1:
%print "It is certain"
%
%elif answers == 2:
%print "Outlook good"
%
%elif answers == 3:
%print "You may rely on it"
%
%elif answers == 4:
%print "Ask again later"
%
%elif answers == 5:
%print "Concentrate and ask again"
%
%elif answers == 6:
%print "Reply hazy, try again"
%
%elif answers == 7:
%print "My reply is no"
%
%elif answers == 8:
%print "My sources say no"
%\end{lstlisting}
%\end{latin}
